\chapter{Conclusion and outlook} \label{chap:summary}

In this final chapter the studied project and its findings are summarised.
Then, an outlook about further investigations of remaining uncertainties are given.

In this thesis \acrshort{tbl} along curved walls are investigated.
The studies are performed with simulations based on \acrshort{dns} and \acrshort{rans}.
With these simulations effects of convex curvatures on the development of the \acrshort{tbl} are studied.
However, in flows along a curvature a pressure gradient arises naturally induced by the geometry.
This pressure gradient also affects the \acrshort{tbl}.
In order to study the effect of the curvature separately, the effect of the pressure gradient needs to be known.
Therefore, a manipulation of the pressure field is required.
In this work we manipulate the pressure such that a \acrshort{zpg} develops.
This is a favourable pressure distribution, since \acrshort{tbl} along flat plate walls also have a \acrshort{zpg} as long as no pressure gradient is imposed.
Therefore, \acrshort{tbl} under \acrshort{zpg} conditions are well-studied and provide an established reference setup.
In order to create a \acrshort{zpg} in a curved simulation domain an optimisation is employed.
The optimisation algorithm introduced by \citet{morita_applying_2022} performs a shape optimisation in a \acrshort{rans} simulation.
The upper wall is defined by several \benennung{control points} which are connected through a spline.
The position of these \benennung{control points} is optimised such that the Clauser parameter $\beta$ at the lower wall fits its target value.
The Clauser parameter corresponds to a non-dimensionalisation of the pressure gradient.
Through the adaptation of the \benennung{control points} the shape of the domain changes, which implicitly affects the pressure distribution.
When the optimisation is converged, the result is imposed on a domain for a \acrshort{dns}.
To do this, the velocities along a specific wall-normal height are extracted.
This height corresponds to the height of the upper wall of the domain used for \acrshort{dns}.
There the extracted velocities are imposed as a \benennung{Dirichlet} \acrshort{bc}.
With this optimised \acrshort{bc} three different curved setups are simulated.
Two setups have a radius of $100\delta_0$ with curvatures of $30$° and $90$°.
The third setup has a curvature of $500 \delta_0$ and $18$° curvature.
All setups have a straight inflow region before the curvature starts to ensure a proper development of the \acrshort{tbl}.
The \acrshort{tbl} is generated with an induced tripping of a laminar Blasius boundary layer near the inflow.
At the end of the curvature a straight outlet region in tangential direction is added as a rebalancing zone.
With these setups two simulations for each curved case are performed.
One simulation is performed with the optimised velocity \acrshort{bc} at the upper wall.
As a base configuration, all cases are also simulated with a \benennung{Dirichlet} \acrshort{bc} of $(1,0,0)$ velocity profile at the upper wall.
With these base simulations the effect of the optimisation can be studied.
In addition, the optimisation capability for the different curvature cases is compared.
The results show, that the optimisation capability differs for the \benennung{local} and \benennung{global} trend of the pressure distribution.
On a \benennung{global} level the optimised for all three cases have a more constant pressure distribution in the curvature compared to their base cases.
However, fluctuations in the pressure distribution arise  \benennung{locally} which are also present in the distribution of the Clauser parameter.
The effects of different curvatures are studied through a comparison of case R100deg90 and R500deg18.
XXX curvature effects 
XXX optimisation for different cases


In this thesis, the workflow, introduced in \citet{schlatter_leveraging_2025}, to generate a curved \acrshort{tbl} domain with \acrshort{zpg} was tested.
The results presented in \autoref{chap:results} show, that the application of the optimisation algorithm presented by \citet{morita_applying_2022} was possible in all three curvature configurations.
The \benennung{global} trend of the pressure field was successfully manipulated towards a constant value in the curved segment.
However, \benennung{local} fluctuations in the pressure distribution were detected.
The positioning of the fluctuations suggest a connection with the position of the \benennung{control points} used in the optimisation.
The studies performed by \citet{ganagdhar_bayesian_2025} about the optimisation capability captured a case with significantly lower curvature compared to the boundary layer thickness.
Therefore, only rare knowledge exists about the capability of the optimisation algorithm for the here tested geometric properties.
In order to eliminate the \benennung{local} fluctuations in the pressure distribution, a detailed study of the optimisation algorithm with various settings and geometries is required.
In addition, the number and positioning of the \benennung{control points} should be investigated to obtain an optimal relation between a smooth $\beta$ distribution and a reasonable computational effort for the optimisation.


In the \acrshort{dns} of the curved \acrshort{tbl} we visibly detected fluctuations in between the instantaneous pressure fields.
Because of their characteristics, analysed in \autoref{sec:pressFluc}, we assume that they are non-physical artefact caused by the outer constrains of the domain.
Similar pressure modes were detected by \citet{maday_spectral_1989} and \citet{sabbah_filtering_2001} who in introduced methods to eliminate them.
We performed various tests with different configurations to localise potential causes for the fluctuations.
However, with these tests we were not able to clearly identify the causes of the fluctuations.
Here, further tests are necessary to find the original source of these fluctuations.
Nevertheless, we verified the reliability of the statistical results to ensure that the statistics unaffected by the fluctuations.
To do this, we extracted data in the straight section and compared them to the earlier performed simulations of a plane \acrshort{tbl} as well as with reference data of plane \acrshort{tbl} simulations.
This verification is presented in \autoref{sec:planeTBLVerif}, where it can be seen that the statistical quantities are undisturbed and hence similar to the reference data.
In order to obtain a better understanding of the fluctuations further tests need to be done.
This also helps for future improvements of the domain setups or the comparisons of different setups.


